\section{Limitations}
\label{sec:limitations}

\subsection{Cohort size and single-seed fragility}
The primary UCSD Stress cohort used throughout
Sec.~\ref{sec:results_audit}--\ref{sec:results_ceiling} consists of 70
recordings from 17 subjects, with a 14:56 positive:negative class
split. At this scale, single-seed balanced accuracies are known to
swing by $5$--$10$\,pp on the same pipeline with no code changes
\citep{smalln2025biorxiv}. All headline numbers are reported as
three-seed mean $\pm$ sample standard deviation; seed counts are
held constant across the trial-level / subject-level / class-balance
/ architecture / HP arms. A reader interpreting a given three-seed
point estimate should weight it against the $\pm 5$\,pp (typical) to
$\pm 10$\,pp (worst observed) envelope; the relative claims, not the
absolute BA digits, carry the weight of the argument. A consequence
for the permutation-null analysis is that our comparison budget is
ten permutation seeds per FM, below the hundreds that would drive
$p$-floor well below $0.01$. For CBraMod and REVE this means the
null test bottoms out at $p \approx 0.1$. The SDL diagnostic does not
hinge on these two cells --- the paired comparison carries the
thesis's weight --- but a larger permutation pool would sharpen the
audit claim and is a natural replication target.

\subsection{Single anchor on the rescue arm}
Our controlled paired comparison anchors one rescue dataset (EEGMAT,
Klimesch-type alpha desynchronisation) against one bounded dataset
(UCSD Stress longitudinal DSS, absent literature anchor). A natural
follow-up would use multiple anchored contrasts --- at minimum
EEGMAT, ADFTD, and perhaps a motor-imagery beta-ERD task --- and
multiple bounded contrasts, to map a contrast-strength gradient
rather than a binary. The observation in
Sec.~\ref{sec:results_fm_adds} that ADFTD produces a similar
classical-to-frozen-to-FT cascade to EEGMAT is a preliminary step in
this direction, but ADFTD's contrast-strength anchor is cross-subject
rather than within-subject, and is therefore not directly comparable
to the paired design.

\subsection{Architecture panel breadth and per-model HP selection}
The architecture panel in Sec.~\ref{sec:results_ceiling} includes
seven from-scratch and pretrained models spanning five orders of
magnitude in parameter count but omits two classes worth flagging:
graph-based EEG models over sensor adjacency, and state-space /
Mamba-style models. Our panel is representative of the convolutional
and transformer families currently dominant in EEG-FM literature, but
a reader who suspects an inductive bias specific to these omitted
families might escape the ceiling has a legitimate concern the panel
does not directly address.

Separately, the three FM fine-tuning results use per-model best-HP
selection over a small $3 \times 2 \times 3$ grid independently per
model. This is an argmax over HP selection, not a factorial
cross-model comparison. The per-model headline numbers
(LaBraM $0.524$, CBraMod $0.548$, REVE $0.577$) should be read as
``each FM at its own best tested HP,'' not as a controlled ranking
between architectures. The architecture-independent-ceiling claim in
Sec.~\ref{sec:results_ceiling} does not rely on cross-FM ranking; it
relies only on every tested architecture at its own best HP falling
inside the same $5$\,pp window.

\subsection{Reproducibility}
All code, trained checkpoints, variance-decomposition outputs,
permutation-null logs, per-seed training curves, the pre/post LaBraM
diagnostic, the class-balanced classical audit, and the contrast-
strength anchoring protocol specification will be released at the
project repository upon publication. We invite independent groups to
(i) reproduce the trial-vs-subject CV decomposition on their own
Stress pipelines, (ii) replicate the paired EEGMAT vs Stress
longitudinal comparison under their own FM weights and HP budgets,
and (iii) apply the pre-benchmark contrast-anchoring protocol to
their own clinical cohorts and report whether SDL's prediction
holds.
